\documentclass[a4paper,12pt]{article}

\usepackage[
  top=2.5cm,
  bottom=2.2cm,
  left=3.0cm,
  right=2.2cm
]{geometry}
\usepackage{fontspec}
\usepackage{setspace}
\usepackage{url}
\usepackage[
backend=biber,
style=authoryear,
maxcitenames=2,
maxbibnames=99
]{biblatex}
\addbibresource{apitan_2020.bib}
\pagestyle{empty}
\urlstyle{same}

\XeTeXlinebreaklocale "th_TH"
\XeTeXlinebreakskip = 0pt plus 1pt

\setmainfont[
  Path=./,
  UprightFont={THSarabunNew.ttf},
  BoldFont={THSarabunNew_Bold.ttf},
  ItalicFont={THSarabunNew_Italic.ttf},
  BoldItalicFont={THSarabunNew_BoldItalic.ttf}
]{THSarabunNew}

\newcommand{\fieldline}[2]{%
  \noindent\hangindent=1.55cm\hangafter=1%
  \makebox[1.55cm][l]{#1}#2\par
}

\begin{document}
	\fontsize{16pt}{20pt}\selectfont
	\setlength{\parindent}{1.27cm}
	\setlength{\parskip}{5pt}
	\sloppy

\begin{flushright}
วันที่ ๓๐ มิถุนายน ๒๕๖๙
\end{flushright}

\vspace{0.35cm}

\fieldline{เรื่อง}{ขอความอนุเคราะห์ข้อมูลราคาสินค้าและบริการรายหมวดย่อยเพื่อใช้ในการศึกษาวิจัย}
\fieldline{เรียน}{ผู้อำนวยการสำนักงานนโยบายและยุทธศาสตร์การค้า}
\fieldline{อ้างอิง}{\fullcite{APAITAN2020117}}

\vspace{0.2cm}

ผม ธนกฤต เมธเศรษฐ นักศึกษาระดับปริญญาเอก คณะเศรษฐศาสตร์ University of Cambridge กำลังดำเนินการศึกษาวิจัยเกี่ยวกับการประมาณแนวโน้มเงินเฟ้อของประเทศไทย โดยมีวัตถุประสงค์เพื่อพัฒนาแบบจำลองสำหรับประมาณแนวโน้มเงินเฟ้อจากข้อมูลราคาสินค้าและบริการรายหมวดย่อย ภายใต้บริบทที่ราคาบางรายการมีลักษณะไม่เปลี่ยนแปลงเป็นเวลานาน (price staleness) ซึ่งอาจส่งผลต่อการวัดแนวโน้มเงินเฟ้อที่แท้จริง ทั้งนี้ งานวิจัยดังกล่าวคาดว่าจะพัฒนาเป็นส่วนหนึ่งของวิทยานิพนธ์ระดับปริญญาเอกของผม

ในการศึกษานี้ ผมมีความประสงค์จะขอความอนุเคราะห์ \textbf{ข้อมูลรายเดือนของราคาสินค้าและบริการรายหมวดย่อยของประเทศไทยที่ใช้ในการจัดทำดัชนีราคาผู้บริโภค} พร้อมทั้งน้ำหนักที่ใช้ในการคำนวณดัชนี หมวดหมู่สินค้า ช่วงเวลาของข้อมูล ผมขอเรียนเพิ่มเติมว่า งานศึกษาของ \textcite{APAITAN2020117} ได้ใช้ข้อมูลราคาสินค้าและบริการรายย่อยของประเทศไทยจากกระทรวงพาณิชย์ ซึ่งเป็นข้อมูลรายเดือน ครอบคลุมสินค้าจำนวนมากใน 77 จังหวัด งานดังกล่าวสะท้อนให้เห็นถึงประโยชน์ของข้อมูลรายหมวดย่อยในการทำความเข้าใจพลวัตเงินเฟ้อของไทยอย่างลึกซึ้ง ผมจึงมีความประสงค์จะขอข้อมูลในลักษณะเดียวกัน หรือข้อมูลในระดับที่กระทรวงพาณิชย์สามารถอนุเคราะห์ได้ เพื่อใช้ต่อยอดในการศึกษาด้านการประมาณแนวโน้มเงินเฟ้อภายใต้ปัญหาราคาที่ไม่เปลี่ยนแปลงเป็นเวลานาน

ข้อมูลที่ได้รับจะถูกนำไปใช้เพื่อวัตถุประสงค์ทางวิชาการเท่านั้น และผมยินดีปฏิบัติตามเงื่อนไขการใช้ข้อมูล การรักษาความลับ หรือข้อตกลงอื่น ๆ ที่ทางกระทรวงพาณิชย์ กำหนดอย่างเคร่งครัด หากจำเป็น ผมยินดีจัดส่งหนังสือรับรองจากสถาบัน หรือเอกสารประกอบเพิ่มเติมเพื่อพิจารณา

จึงเรียนมาเพื่อพิจารณา

\vspace{1.5cm}

\begin{flushright}
\begin{minipage}{9cm}
\centering
%ขอแสดงความนับถือ
%
%\vspace{1.25cm}

(ธนกฤต เมธเศรษฐ)\\
นักศึกษาระดับปริญญาเอก University of Cambridge\\
tm921@cam.ac.uk\\
(+66) 0954989455
\end{minipage}
\end{flushright}

\end{document}
